% Guia do professor — Capítulo 10: Forças e Ventos
\section{Capítulo 10 --- Forças Atmosféricas e Ventos}

\subsection*{Visão geral e ritmo}

A porta de entrada da dinâmica: Coriolis, equilíbrio geostrófico, vento
gradiente, atrito/Ekman e continuidade — terminando na frase que organiza a
sinótica: \emph{baixas convergem e chovem; altas divergem e abrem}.
\textbf{Ritmo: 2 aulas de 2\,h.} Para físicos, é o capítulo mais
confortável do curso — aproveite para ir fundo nas deduções.

\begin{itemize}
  \item \textbf{Aula 1} — forças, Coriolis (com a desmistificação da pia),
        círculos inerciais, dedução geostrófica completa, Buys-Ballot,
        régua na carta. Casos~1--2 (15\,min).
  \item \textbf{Aula 2} — vento gradiente (quadrática no quadro), 3 forças
        com atrito e o desenho convergência/divergência, Ekman,
        continuidade e $w$. Casos~3--4 (15\,min).
\end{itemize}

\subsection*{Objetivos de aprendizagem}

\begin{enumerate}
  \item escrever a equação do movimento horizontal e a escala de cada
        termo;
  \item deduzir o vento geostrófico e usá-lo como régua de carta
        (Buys-Ballot no HS);
  \item explicar sub/supergeostrofia em baixas/altas e o regime
        ciclostrófico;
  \item desenhar o equilíbrio de 3 forças e concluir
        convergência$\to$subida em baixas;
  \item derivar/verificar a espiral de Ekman e estimar $w$ por
        continuidade.
\end{enumerate}

\subsection*{Pré-requisitos a reativar}

Referenciais não inerciais e força centrífuga/Coriolis (mecânica);
gradiente e divergência (cálculo vetorial); hidrostática (Cap.~1) para a
versão em coordenadas de pressão (mencionar, não exigir).

\subsection*{Armadilhas conceituais frequentes}

\begin{itemize}
  \item \textbf{A pia}: Coriolis em escala de pia é $10^{-7}$ da força
        real; medir com o T1 (escala de tempo).
  \item \textbf{HS × HN}: todo desenho de livro americano gira ao
        contrário aqui; treinar Buys-Ballot no HS até virar reflexo.
  \item \textbf{``Coriolis faz o vento''}: Coriolis só \emph{desvia}; quem
        empurra é a PGF.
  \item \textbf{Geostrofia como lei}: é equilíbrio assintótico; perto do
        equador, em curvas fechadas e na CL ele falha — e as falhas são o
        conteúdo dos Caps.~11--16.
  \item \textbf{Sinal de $f$}: erros de sinal no HS são epidêmicos em
        listas; cobrar o desenho junto com a conta.
\end{itemize}

\subsection*{Demonstração de baixo custo}

Globo girante (ou cadeira giratória) + bolinha rolada: o desvio aparente
ao filmar do referencial girante. Vídeo clássico do carrossel do MIT
fecha a discussão da pia.

\subsection*{Problemas sugeridos do livro (exercícios do Cap.~10)}

Aquecimento: $f$ em várias latitudes; $U_g$ de cartas dadas; período
inercial. Aprofundamento: vento gradiente em baixas/altas; ângulo de
cruzamento com atrito. Síntese: ``por que o centro de uma alta é sempre de
ventos fracos?'' (T3 responde).
\emph{Confira a numeração na sua cópia (v1.02b).}

\subsection*{Uso do notebook \texttt{cap10\_ventos.ipynb}}

\begin{center}
\begin{tabular}{lll}
\hline
Caso & Conteúdo & Momento sugerido \\
\hline
1 & Círculos inerciais e cicloides & Aula 1, após Coriolis \\
2 & Geostrófico × gradiente na baixa sintética & Aula 1/2, após dedução \\
3 & Ekman: analítica × numérica & Aula 2, após atrito \\
4 & Convergência $\to$ mapa de $w$ & fim da Aula 2 \\
\hline
\end{tabular}
\end{center}

O Caso~4 é o fecho conceitual: o mapa de $w$ da baixa sintética É o mapa
de nuvens que os alunos verão nas imagens de satélite do Cap.~13.

\subsection*{Exercícios complementares e soluções}

Nas notas: T1--T5 (inerciais, equador, teto das altas, ângulo de
cruzamento, verificação de Ekman + transporte) e N1--N4. Soluções em
\texttt{cap10\_solucoes.ipynb} (\emph{não distribuir}): T3 e T5 com
\texttt{sympy}. Pares: T1$\leftrightarrow$N1, T5$\leftrightarrow$N3.

\subsection*{Conexões com o resto do curso}

Isóbaras$\to$vento fecha o Cap.~9; convergência$\to w$ abre os
Caps.~12--13; Ekman e cruzamento $\to$ camada limite (Cap.~18) e furacões
(Cap.~16); ciclostrófico $\to$ tornados (Cap.~15); geostrofia em altitude
$\to$ vento térmico (Cap.~11).
